% ============================================================
%  §4.3 Δημιουργία διαγράμματος GGV
% ============================================================

\label{subsec:ggv}

Το διάγραμμα GGV ($G$-$G$-$V$) αποτελεί τον πυρήνα της quasi-static προσομοίωσης γύρου:
ορίζει, για κάθε ταχύτητα $v$, το σύνολο των εφικτών συνδυασμών επιτάχυνσης $(a_x, a_y)$
που το όχημα μπορεί να παράγει χωρίς να ξεπεράσει τα όρια πρόσφυσης. Στη συνέχεια, ο lap
simulator χρησιμοποιεί αυτό το envelope για να βρει, σε κάθε σημείο πίστας, τη μέγιστη
επιτρεπτή ταχύτητα.

\subsubsection{Δομή του διαγράμματος}

Για κάθε τιμή ταχύτητας $v_j$ σε ένα πλέγμα $v_{\min} \leq v \leq v_{\max}$,
κατασκευάζεται μια «φέτα» του GGV: ένα σύνολο σημείων $(a_{x,k}, a_{y,\max,k})$ στο
επίπεδο $G$-$G$. Κάθε σημείο αντιστοιχεί σε ένα συγκεκριμένο επίπεδο διαμήκους επιτάχυνσης
$a_{x,k}$, και ο αλγόριθμος υπολογίζει τη μέγιστη πλευρική επιτάχυνση που είναι εφικτή υπό
αυτό το $a_x$ -- αυτό είναι η \emph{friction ellipse} υπό φορτίο.

Η μαθηματική σχέση που συνδέει $a_y$ με $a_x$ δεν είναι αναλυτική: εξαρτάται από τη
μη-γραμμική ανταπόκριση ελαστικών, την κατανομή φορτίου, και τη θερμοκρασία. Γι' αυτό η
επίλυση γίνεται αριθμητικά με δυαδική αναζήτηση επί της ακτίνας καμπυλότητας $r$.

\subsubsection{Αλγόριθμος δυαδικής αναζήτησης}

Για δεδομένο $(v, a_x)$, η μέγιστη $a_y$ αντιστοιχεί στην ελάχιστη ακτίνα καμπυλότητας
$r_{\min}$ για την οποία ακόμα επιτυγχάνεται ισορροπία δυνάμεων. Η δυαδική αναζήτηση
(bisection) λειτουργεί:

\begin{enumerate}
    \item Ορίζεται εύρος $[r_{\rm lo}, r_{\rm hi}]$ -- από μία πολύ μεγάλη ακτίνα (ευθεία, χωρίς πλευρική επιτάχυνση) έως μία πολύ μικρή (πολύ μεγάλη καμπύλη, οπωσδήποτε εκτός ορίων).
    \item Δοκιμάζεται $r_{\rm mid} = (r_{\rm lo} + r_{\rm hi})/2$: αν ο solver συγκλίνει και τα ελαστικά δεν κορεννύονται, το $r_{\rm lo}$ εκ νέου μικρύνεται· αλλιώς το $r_{\rm hi}$ μεγαλώνει.
    \item Το κριτήριο σύγκλισης είναι $r_{\rm hi} - r_{\rm lo} < \varepsilon$. Τότε $r_{\min} \approx r_{\rm lo}$ και:
\end{enumerate}

\begin{equation}
    a_{y,\max} = \frac{v^2}{r_{\min}}
    \label{eq:ay_from_r}
\end{equation}

Κάθε κλήση του solver εκτελεί πλήρη επίλυση 2-DOF (ίδια δομή με τη
§\ref{subsec:slip_estimation}), συμπεριλαμβανομένης της κατανομής φορτίου, των ολισθήσεων
και των δυνάμεων ελαστικού. Επομένως ο υπολογισμός ενός πλήρους GGV (π.χ. 20 τιμές $v$ ×
15 τιμές $a_x$ × ~12 iterations bisection × 4 evaluations/iteration) αντιστοιχεί σε
χιλιάδες κλήσεις solver.

\subsubsection{Θερμοκρασιακή παράμετρος}

Κάθε GGV παράγεται υπό συγκεκριμένη θερμοκρασία ελαστικών $T_{\rm ref}$: η συνάρτηση
$g(T_{\rm ref})$ κλιμακώνει τον συντελεστή τριβής $\mu$ σε ολόκληρο τον υπολογισμό. Στην
πράξη παράγεται οικογένεια GGV (διαφορετικές $T_{\rm ref}$) και ο lap simulator επεμβαίνει
με παρεμβολή -- αντιστοιχεί σε θερμοκρασιακά μεταβαλλόμενο envelope κατά τη διάρκεια γύρου.

% TODO: Αξίζει να δείξεις τη σύγκριση GGV ψυχρού vs. θερμού ελαστικού.
% Η οπτική διαφορά του envelope είναι η πιο εντυπωσιακή απόδειξη
% ότι η θερμοκρασία αλλάζει ουσιαστικά τη δυναμική του οχήματος.

% ----- TikZ: GGV pipeline flowchart -----
\begin{figure}[htbp]
\centering
\begin{tikzpicture}[
    node distance=0.50cm and 1.0cm,
    block/.style={rectangle,draw=black,fill=blue!8,text width=3.8cm,
                  align=center,rounded corners=4pt,minimum height=0.88cm,font=\small},
    decision/.style={diamond,draw=black,fill=orange!15,text width=2.4cm,
                     align=center,font=\scriptsize,inner sep=1pt},
    io/.style={rectangle,draw=black!60,fill=green!10,text width=3.4cm,
               align=center,rounded corners=2pt,font=\small,minimum height=0.78cm},
    outbox/.style={rectangle,draw=black!60,fill=red!8,text width=3.8cm,
                align=center,rounded corners=2pt,font=\small,minimum height=0.78cm},
    arr/.style={->,thick,gray!70!black}
]
  \node[io] (inp) {$v_j$, $a_{x,k}$, $T_{\rm ref}$\\πλέγμα εισόδου};

  \node[block,below=of inp] (init)
    {Αρχικοποίηση\\$r_{\rm lo}$, $r_{\rm hi}$};
  \draw[arr] (inp)--(init);

  \node[block,below=of init] (mid)
    {$r_{\rm mid} = (r_{\rm lo}+r_{\rm hi})/2$\\$a_y = v^2 / r_{\rm mid}$};
  \draw[arr] (init)--(mid);

  \node[block,below=of mid] (solver)
    {2-DOF solver\\$\delta$, $\beta$ για $(v,a_x,a_y)$\\+ tyre forces};
  \draw[arr] (mid)--(solver);

  \node[decision,below=0.6cm of solver] (feas)
    {feasible?};
  \draw[arr] (solver)--(feas);

  % yes branch: tighten r_lo
  \node[font=\scriptsize,left=1.2cm of feas] (tighten) {$r_{\rm lo} \leftarrow r_{\rm mid}$};
  \draw[arr] (feas.west)--node[above,font=\scriptsize]{ναι}(tighten.east);
  \draw[arr,rounded corners=4pt] (tighten.north)|-(mid.west);

  % no branch: widen r_hi
  \node[font=\scriptsize,right=1.2cm of feas] (widen) {$r_{\rm hi} \leftarrow r_{\rm mid}$};
  \draw[arr] (feas.east)--node[above,font=\scriptsize]{όχι}(widen.west);
  \draw[arr,rounded corners=4pt] (widen.north)|-(mid.east);

  \node[decision,below=0.6cm of feas] (conv) {$r_{\rm hi}-r_{\rm lo}<\varepsilon$?};
  \draw[arr] (feas)--(conv);

  \node[block,below=0.6cm of conv] (store)
    {$a_{y,\max}(v_j,a_{x,k}) = v_j^2/r_{\rm lo}$\\αποθήκευση};
  \draw[arr] (conv)--node[right,font=\scriptsize]{ναι}(store);

  \draw[arr,dashed,rounded corners=4pt]
    (conv.east)--++(0.5,0)
    node[right,font=\scriptsize]{όχι $\Rightarrow$ επανάληψη};

  \node[outbox,below=of store] (outp)
    {\textbf{Έξοδος}: GGV array\\$a_{y,\max}(v,a_x,T_{\rm ref})$};
  \draw[arr] (store)--(outp);
\end{tikzpicture}
\caption{Pipeline υπολογισμού διαγράμματος GGV μέσω δυαδικής αναζήτησης. Για κάθε ζεύγος $(v_j, a_{x,k})$ βρίσκεται η ελάχιστη εφικτή ακτίνα καμπυλότητας, άρα η μέγιστη $a_y$.}
\label{fig:ggv_flow}
\end{figure}
% ----- τέλος flowchart -----

